\documentclass[10pt]{article}

\usepackage[utf8]{inputenc}

\usepackage[T1]{fontenc}

\usepackage{amsmath}

\usepackage{amsfonts}

\usepackage{amssymb}

\usepackage[version=4]{mhchem}

\usepackage{stmaryrd}

\usepackage{mathrsfs}



\begin{document}

Пусть $V$ - регулярно определимое многообразие У. а., заданное множеством $\left\{f_{i} \mid i \in I\right\}, I \neq \varnothing$, функциональных символов и множеством $\Sigma$ тождеств. Каждому символу $f_{i}$ сопоставляется элемент $a_{i}$, а для каждого тождества вида $\mathrm{I}_{1}$ из $\Sigma$ выписывается определяющее соотношение



\[

a_{i_{1}} \ldots a_{i_{k}}=a_{j_{1}} \ldots a_{i_{l}} .

\]



Пусть $P$ - полугруппа с порождающими $a_{i}, i \in I$, и выписанными определяющими соотношениями, а $P^{1}-$ полугруппа $P$ с внешне присоединенной единицей $e$. Для каждого тождества вида $\mathrm{I}_{2}$ из $\Sigma$ (если такие имеются) выписывают определяющее соотношение $a_{i_{1}} \ldots a_{i_{k}}=e$. Полугруппу $P_{V}$, получаемую из $P^{1}$ присоединением всех таких определяющих соотношений, и считают соответствующей многообразию $V$. Она во многом характеризует әто многообразие. Если $\Sigma$ содержит лишь тождества вида $\mathrm{I}_{1}$, то можно ограничиться построением лишь полугруппы $P$. Определяя в $P_{V}$ унарные операции $f_{i}(x)=x a_{i}$, получают У. а. $\left\langle P_{y},\left\{f_{i} \mid i \in I\right\}\right\rangle$, к-рая является $V$-свободной алгеброй ранга 1. Группа всех автоморфизмов У. а. $\left\langle P_{V},\left\{f_{i} \mid i \in I\right\}\right\rangle$ изоморфна группе $P_{V}^{*}$ обратимых элементов полугруппы $P_{V}$.



Лum.: [1] М а л ь ц е в А. И., Алгебраические системы, М., 1970; [2] Б и р к г о Ф Г., Б а р т и Т., Современная прикладная алгебра, пер. с англ., М., 1976; [3] С̈ м и р н о в Д. М., "Алгебра и логика", 1976, т. 15, № 3, с. 331-42; [4] е г о же, там же, 1978, т. 17, № 4, c 468-77, [5] J on s s o n B., Topics in universal algebra, B.- Hdlb - N. Y., 1972. Д. М. Смирнов.



УНИВЕРСАЛЬНАЯ АЛГЕБРА - алгебраическая система с пустым множеством отношений. У. а. часто называют просто а л г е б р о й. Для У. а. справедлива те о рем а о гом о мо рфизме: если $\varphi-$ гомоморфизм У. а. $A$ на У. а. $B$ и $\theta-$ ядерная конгруэнция гомоморфизма $\varphi$, то $B$ изоморфна факторалгебре $A / \theta$. Всякая У. а. разлагается в подпрямое произведение подпрямо неразложимых У. а.



Если к основным операциям алгебры $A$ присоединить все производные операции, то возникает У. а. $\tilde{A}$ большей сигнатуры. Равенство $\tilde{A}=\tilde{B}$ возможно и при $A \neq B$, что приводит $к$ понятию рациональной әквивалентности У. а. (см. Универсальных алгебр многообразие).



С каждой У. а. $A$ связаны с о п у т с т в у ю ш и е с т р у к т у р ы: моноид всех эндоморфизмов End $A$, группа всех автоморфизмов Aut $A$, решетки всех подалгебр Sub $A$ и всех конгруэнций Con $A$. Для любых группы $G$ и аләебраических решеток $U$ и $C$ существует такая У. а. $A$, что $G \cong$ Aut $A, U \cong \operatorname{Sub} A$ и $C \cong \operatorname{Con} A$ (см. [12]). Однако при замене Aut $A$ на End $A$ аналогичный результат не имеет места. Такого рода задачи наз. а б с т р а т тым и задачам и р е а ли зац и и. Пример решения к о н к е т но й 3 а д а ч и р е а л и з а ц и и: система подмножеств $U$ множества $A$ совпадает с Sub $A$ для нек-рой У. а. с носителем $A$ тогда и только тогда, когда $U$ замкнута относительно объединения направленных подсистем и произвольных пересечений [11]. Как абстрактную, так и конкретную задачи реализации можно решать и для заданного класса У. а. Исследовались У. а. с теми или иными ограничениями на сопутствующие структуры. Напр., У. а. с дистрибутивной или дедекиндовой решеткой конгруәнций, с двуәлементной решеткой конгруәнций (к о н г р у ә н ц-П р о с т ы е У. а.), с одноәлементной или двуэлементной решеткой подалгебр (п р о с т ы е У. а.), с коммутативным моноидом әндоморфизмов, с одноәлементной группой автоморфизмов (Ж е с т к и е У. а.) и т. П. У. а. с перестановочными конгруәнциями изоморфна прямому произведению конечного числа конгруәнц-шростых алгебр тогда и только тогда, когда решетка ее конгруэнций удовлетворяет условию максимальности, а точная верхняя грань ее минимальных конгруэнций равна наибольшей конгруэнции. У. а. с дистрибутивной решеткой конгруәнций и перестановочными конгруэнциями (а р и ф м е т ич е с к и е У. а.) допускают представление в виде глобального сечения подходящих пучков. Исследовалось, насколько $\mathrm{Y}$. а. определяется той или иной из своих сопутствующих структур. Впрочем, большинство результатов такого рода касается конкретных глассов У. a. ([9] - [12], [15]).



У. а. наз. ф у н к ц и о нально по лн о ї, если всякая операция на ее носителе принадлежит клону, порожденному ее основными операциями и константами. Если отказаться от включения констант, то получается определение п р и м а л ь н о й (или с т р о го ф ункци о на ль н о по лн о й) У. а. Принадлежность к вышеупомянутому клону всех огераций, сохраняемых конгруэнциями, определяет а ф ф и н н о п о л н у ю У. а. Всякая функционально полная У. а. конечна. Поэтому требование конечности часто включают в определение перечисленных классов У. а. (см. $[9],[13]$, [14]).



Формирование теории У. а. началось в $30-40-х$ гг. 20 в., когда были сформулированы основные определения, охарактеризованы многообразия универсальных алгебр и доказана теорема о подпрямых разложениях (см. [7], [8]). Предыстория теории У. а. восходит к прошлому столетию. Активные исследования в этой области начались в кон. 40-х гг., а в СССР - в нач. 50-х гг. (А. Г. Курош, А. И. Мальцев и их ученики). Привлечение методов математич. логики привело к рассмотрению алгебрачческих систем.



Термин «У. а.» употребляется также в смысле «теория универсальных алгебр".



Лum.: [1] Б и $\mathrm{p}_{\mathrm{K}} \mathrm{r}$ о ф Г., Теория решеток, пер. с англ., 2 изд., М., 1983; [2] К о н П., Универсальная алгебра, пер. с англ., М., 1968; [3] К у р о шा А. Г., Лекции по общей алгебре, 1970 учебного года, М , $1974 ;$; [5] М а ль ь ц е в А. И., Алгебраические системы, М., 1970, [6] С к о р н я к о в JI. A., Элементы общей алгебры, М., 1983: [7] В i r k h of f G., „Proc. Cambridge Phil. Soc.», 1935, v. 31, p. 433-54; [8] e r o ж e, "Bull. Amer. Math. Soc.», 1944, v 50, p. 764-68; [9] G r a t z e r G., Universal algebra,"2 ed., N. Y.- 'Hdlb.- B., 1979; [10] H u T a hK a i, "Math. Nachr.", 1969, Bd 42, N i-3, S. 157-71; [11] $\mathrm{J}$ ón s s o n B., Topics in universal algebra, B.- Hdlb. - N. Y., 1972; [12] L a m p e W. A., "Algebra universalis", 1972, v. 2, "Coll. Math. Soc. J. Bolyai», 1982, v. 29, p 583-608, [14] W e rn e r H., Discriminator-algebras, B., 1978; [15] W o l f A., "Mem. Ámer. Math. Soc.», 1974, v. 148, p. 87-93.



УНИВЕРСАЛЬНАЯ ОБЕРТЫВАЮШАЯ АЛГЕБРА а лге б ры Л и $\mathfrak{g}$ на д ком м у т а тивны м к о ль ц о м $k$ с е д и н и ц е й - ассоциативная $k$ алгебра $U(\mathfrak{g})$ с единицей, снабженная отображением $\sigma: \mathfrak{g} \rightarrow U(\mathfrak{g})$, для к-рой выполнены следующие свойства: 1) $\sigma$ является гомоморфизмом алгебр Ли, т. е. $\sigma$ $k$-линейно и $\sigma([x, y])=\sigma(x) \cdot \sigma(y)-\sigma(y) \cdot \sigma(x), x, y \in \mathfrak{g}$; 2) для любой ассоциативной $k$-алгебры $A$ с единицей и всякого такого $k$-линейного отображения $\alpha: \mathfrak{g} \rightarrow A$, тто $\alpha([x, y])=\alpha(x) \alpha(y)-\alpha(y) \alpha(x), x, y \in \mathfrak{g}$, существует единственный гомоморфизм ассоциативных алгебр $\beta: U(\mathfrak{g}) \rightarrow A$, переводящий единицу в единицу, для к-рого $\alpha=\beta$ о б. У. о. а. определяется однозначно с точностью до изоморфизма и всегда существует: если $T(\mathfrak{g})$ - тензорная алгебра $k$-модуля $\mathfrak{g}, I$ - ее двусторонний идеал, порожденный элементами вида $[x, y]$ $-x \otimes y+y \otimes x, x, \quad y \in \mathfrak{g} \quad$ и $\sigma: \mathfrak{g} \rightarrow T(\mathfrak{g}) / I$ - каноническое отображение, то $T(\mathfrak{g}) / I-\mathrm{I}$. о. а. для $\mathfrak{g}$.



Если $k$ нётерово, а модуль $\mathfrak{g}$ конечного порядка, то алгебра $U(\mathfrak{g})$ - нётерова слева и справа. Если $\mathfrak{g}-$ свободный модуль над областью целостности $k$, то $\mathscr{U}(\mathfrak{g})$ не имеет делителей нуля. Для любой конечномерной алгебры Ли $\mathfrak{g}$ над полем $k$ алгебра $U(\mathfrak{g})$ удовлетворяет условию Оре (см. Вложение полугруппы) и тем самым обладает телом частных.





\end{document}